% Part: first-order-logic
% Chapter: axiomatic-deduction
% Section: deduction-theorem

% verification of properties of provability needed for maximally
% consistent sets in the completeness chapter.

\documentclass[../../../include/open-logic-section]{subfiles}

\begin{document}

\iftag{FOL}
      {\olfileid{fol}{axd}{ded}}
      {\olfileid{pl}{axd}{ded}}

\olsection{నిగమన సిద్ధాంతం}

మనం చూసినట్లుగా, స్వీకృతాధారిత వ్యవస్థలో \tetoken{వ్యుత్పత్తులు}{!!{derivation}s} ఇవ్వడం
శ్రమతో కూడిన పని; \tetoken{వ్యుత్పత్తులను}{!!{derivation}s} కనుగొనడం కూడా కష్టం కావచ్చు. పొడవైన
\tetoken{సూత్రాలు}{!!{formula}s} జాబితాలను నిజంగా రాయడానికి బదులు, కొన్ని సరళమైన ఫలితాలను
ఉపయోగించి అలాంటి \tetoken{వ్యుత్పత్తులు}{!!{derivation}s} ఉన్నాయని వాదించడం సాధారణంగా సులభం.
అలాంటి మూడు ఫలితాలను ఇప్పటికే స్థాపించాం:
$!A \in \Gamma$ అని తెలిసినప్పుడు ఎల్లప్పుడూ $\Gamma \Proves !A$
అని చెప్పవచ్చని \olref[ptn]{prop:reflexivity} చెబుతుంది.
$\Gamma \Proves !A$ అయితే $\Gamma \cup \{!B\} \Proves !A$ కూడా అని
\olref[ptn]{prop:monotonicity} చెబుతుంది. ఇంకా, $\Gamma \Proves !A$,
$!A \Proves !B$ అయితే $\Gamma \Proves !B$ అని
\olref[ptn]{prop:transitivity} సూచిస్తుంది. మోడస్ పోనెన్స్‌కు ఒక
``మెటా'' రూపమైన మరొక సరళ ఫలితం ఇదిగో:

\begin{prop}
  \ollabel{prop:mp} $\Gamma \Proves !A$, $\Gamma \Proves !A \lif
  !B$ అయితే, $\Gamma \Proves !B$.
\end{prop}

\begin{proof}
  $\{!A, !A \lif !B\} \Proves !B$ అని మనకు తెలుసు:
  \begin{derivation}
    1. & $!A$ & Hyp.\\
    2. & $!A \lif !B$ & Hyp.\\
    3. & $!B$ & 1, 2, MP
  \end{derivation}
  \olref[ptn]{prop:transitivity} వల్ల, $\Gamma \Proves !B$.
\end{proof}

ఈ సందర్భంలో మనం ఉపయోగించే అత్యంత ముఖ్యమైన ఫలితం నిగమన సిద్ధాంతం:

\begin{thm}[నిగమన సిద్ధాంతం]
\ollabel{thm:deduction-thm} $\Gamma \cup \{!A\} \Proves !B$ అయితే,
అప్పుడు మాత్రమే $\Gamma \Proves !A \lif !B$.
\end{thm}

\begin{proof}
``అయితే'' దిశ తక్షణమే వస్తుంది. $\Gamma \Proves !A \lif !B$ అయితే,
\olref[ptn]{prop:monotonicity} వల్ల $\Gamma \cup \{!A\} \Proves !A
\lif !B$ కూడా. \olref[ptn]{prop:reflexivity} వల్ల $\Gamma \cup
\{!A\} \Proves !A$ కూడా. కాబట్టి \olref{prop:mp} వల్ల $\Gamma \cup
\{!A\} \Proves !B$.

``అప్పుడే'' దిశ కోసం, $\Gamma \cup \{!A\}$ నుంచి $!B$కు గల
\tetoken{వ్యుత్పత్తి}{!!{derivation}} పొడవుపై ఆగమనం చేస్తాం.

ఆగమన ప్రాతిపదికలో, పొడవు~$1$ ఉన్న ప్రతి \tetoken{వ్యుత్పత్తికి}{!!{derivation}} ఈ వాదనను
నిరూపిస్తాం. $\Gamma \cup \{!A\}$ నుంచి~$!B$కు పొడవు~$1$ ఉన్న
\tetoken{ఒక వ్యుత్పత్తిలో}{!!^a{derivation}} $!B$ ఒక్కటే ఉంటుంది; అది సరైనదైతే,
$!B \in \Gamma \cup \{!A\}$ లేదా ఒక స్వీకృతం. $!B \in \Gamma$
అయినా, $!B$ ఒక స్వీకృతమైనా, $\Gamma \Proves !B$. అంతేకాక,
\olref[prp]{ax:lif1} వల్ల $\Gamma \Proves !B \lif (!A \lif !B)$;
అప్పుడు \olref{prop:mp} వల్ల $\Gamma \Proves !A \lif !B$. $!B \in
\{!A\}$ అయితే, $\Gamma \Proves !A \lif !B$; ఎందుకంటే చివరి
\tetoken{వాక్యం}{!!{sentence}}~$!A \lif !B$ అనేది $!A \lif !A$తో సమానం, దాన్ని
\olref[pro]{ex:identity}లో మనం \tetoken{నిగమనం}{!!{derive}} చేశాం.
\sourcecorrection{OLTEAXD-002}{ఆంగ్ల మూలం ఆగమన ప్రాతిపదికలో
``$!B$ is either $\in \Gamma \cup \{!A\}$'' అని వ్యాకరణపరంగానూ,
గణితపరంగానూ వికృతంగా రాసింది. ఉద్దేశించిన సభ్యత్వ వాక్యాన్ని
$!B \in \Gamma \cup \{!A\}$గా సరిచేశాం.}

ఆగమన దశలో, $\Gamma \cup \{!A\}$ నుంచి~$!B$కు గల \tetoken{ఒక వ్యుత్పత్తి}{!!a{derivation}},
మోడస్ పోనెన్స్ ద్వారా సమర్థించబడిన~$!B$ దశతో ముగుస్తుందని అనుకుందాం.
(అది మోడస్ పోనెన్స్ ద్వారా సమర్థించబడకపోతే, $!B \in \Gamma$,
$!B \ident !A$, లేదా $!B$ ఒక స్వీకృతం; అప్పుడు ఆగమన ప్రాతిపదికలోని
అదే వాదన వర్తిస్తుంది.) అప్పుడు, ఏదో \tetoken{సూత్రం}{!!{formula}}~$!C$కు,
ఆ \tetoken{వ్యుత్పత్తిలోని}{!!{derivation}} ముందరి దశలలో $!C \lif !B$, $!C$ ఉంటాయి; అంటే,
$\Gamma \cup \{!A\} \Proves !C \lif !B$, $\Gamma \cup \{!A\}
\Proves !C$. వాటి \tetoken{వ్యుత్పత్తులు}{!!{derivation}s} పొట్టివి కనుక, ఆగమన పరికల్పన
వాటికి వర్తిస్తుంది. అందువల్ల ఈ రెండూ మనకు లభిస్తాయి:
\begin{align*}
  & \Gamma \Proves !A \lif (!C \lif !B); \\
  & \Gamma \Proves !A \lif !C.
\end{align*}
అయితే, \olref[prp]{ax:lif2} వల్ల
\[
\Gamma \Proves (!A \lif (!C \lif !B)) \lif
((!A\lif !C)  \lif (!A \lif !B)),
\]
కూడా; \olref{prop:mp}ను రెండుసార్లు వర్తింపజేస్తే, కావలసినట్లుగా
$\Gamma \Proves !A \lif !B$ వస్తుంది.
\end{proof}

నిగమన సిద్ధాంతం చెల్లేలా చేయడానికే \olref[prp]{ax:lif1},
\olref[prp]{ax:lif2}లను సరిగ్గా ఆ రూపాల్లో ఎంచుకున్నామని గమనించండి.

కిందివి \tetoken{వ్యుత్పాద్యత}{!!{derivability}} గురించిన కొన్ని ఉపయోగకరమైన వాస్తవాలు; వాటిని
సాధనలుగా వదిలేస్తున్నాం.

\begin{prop}
  \ollabel{prop:derivfacts}
\begin{enumerate}
\item $\Proves (!A \lif !B) \lif ((!B \lif !C)
  \lif (!A \lif !C))$; \ollabel{derivfacts:a}
\item $\Gamma \cup \{ \lnot !A\}
  \Proves \lnot !B$ అయితే, $\Gamma \cup \{ !B\} \Proves
  !A$ (విపర్యయనం); \ollabel{derivfacts:b}
\item $\{ !A, \lnot!A\} \Proves
    !B$ (అసత్యం నుంచి ఏదైనా, విస్ఫోటనం); \ollabel{derivfacts:c}
\item $\{ \lnot\lnot!A\} \Proves
  !A$ (ద్వినిషేధ తొలగింపు);\ollabel{derivfacts:d}
\item $\Gamma \Proves \lnot\lnot!A$ అయితే, $\Gamma \Proves
  !A$;\ollabel{derivfacts:e}
\end{enumerate}
\end{prop}
\sourcecorrection{OLTEAXD-003}{ఆంగ్ల మూలంలో మొదటి నిరూప్యతా వాస్తవపు
సూత్రం చివర ఒక కుడి కుండలీకరణ చిహ్నాన్ని వదిలింది. అంతర్గత,
బాహ్య నిబంధనల రెండింటినీ మూసేలా $(!A \lif !C))$గా సరిచేశాం.}

\tagprob{FOL}
\begin{prob}
  \olref[fol][axd][ded]{prop:derivfacts}ను నిరూపించండి.
\end{prob}
\tagendprob

\tagprob{notFOL}
\begin{prob}
  \olref[pl][axd][ded]{prop:derivfacts}ను నిరూపించండి.
\end{prob}
\tagendprob

\end{document}
